
\renewcommand{\Folder}{2025_RecorridoUPV}
\renewcommand{\EntradaBibtex}{2025_RecorridoUPV}

\begin{frame}{\citetitle{\EntradaBibtex} \footnotemark[1] (1)}
\begin{columns}
	\begin{column}{0.60\textwidth}
		\begin{itemize}
			\item Primer aplicación de realidad virtual solicitada a estudiantes en la UPV. 
			\item Ha pasado por 4 generaciones hasta la versión actual.
            \item Proyecto destacado ya que se utiliza para ferias vocaciones y eventos de promoción.
		\end{itemize}
	\end{column}
	\begin{column}{0.40\textwidth}
		\begin{center}
			\includegraphics[width=0.98\textwidth]{\Folder/figs/img1}
		\end{center}
		 \begin{center}
		     \includegraphics[width=0.98\textwidth]{\Folder/figs/img2}
		 \end{center}

	\end{column}
	
\end{columns}

%\end{block} 
\footnotetext[1]{\fullcite{\EntradaBibtex}}
\end{frame}


\begin{frame}{\citetitle{\EntradaBibtex}  (2)}
\begin{columns}
	\begin{column}{0.48\textwidth}
		\begin{itemize}
			\item Se modelan 4 edificios de la UPV (actualmente hay 6). 
			\item Con el control de XBox es posible moverse en el entorno virtual.
			\item Los sensores del smartphone actualizan la vista en base a la posición de la cabeza del usuario.
		\end{itemize}


	\end{column}
	\begin{column}{0.48\textwidth}
		\begin{center}
			\includegraphics[width=0.98\textwidth]{\Folder/figs/img3}
		\end{center}
		 \begin{center}
		     \includegraphics[width=0.98\textwidth]{\Folder/figs/img4}
		 \end{center}

	\end{column}
	
\end{columns}

%\end{block} 
%\footnotetext[1]{\fullcite{\EntradaBibtex}}
\end{frame}

